\documentclass[12pt]{article}
\usepackage{amsmath, amssymb, amsthm}
\usepackage{mathtools}
\usepackage{tensor}
\usepackage{geometry}
\usepackage{booktabs}
\usepackage{array}
\usepackage{xcolor}
\usepackage{hyperref}
\usepackage{fancyhdr}
\usepackage{slashed}
\usepackage{mdframed}
\geometry{margin=1.1in, headheight=15pt}

\definecolor{cadblue}{RGB}{30,90,200}
\definecolor{verifygreen}{RGB}{20,140,60}

\newcommand{\cdbexpr}[1]{\textcolor{cadblue}{$#1$}}
\newcommand{\verified}[1]{\textcolor{verifygreen}{\checkmark\; #1}}
\newcommand{\half}{\tfrac{1}{2}}
\DeclareMathOperator\Tr{Tr}
\newtheorem{theorem}{Theorem}

\pagestyle{fancy}
\fancyhf{}
\lhead{Srednicki QFT --- Chapter 42}
\rhead{Electrodynamics in Coulomb Gauge}
\cfoot{\thepage}

\title{\textbf{Srednicki QFT: Chapter 42}\\[6pt]
       \large Electrodynamics in Coulomb Gauge\\[4pt]
       \normalsize Gauge fixing, photon propagator, Coulomb interaction,\\
       and the path to spinor-helicity}
\author{Cadabra2 + NumPy verification}
\date{}

\begin{document}
\maketitle
\tableofcontents
\newpage

\begin{mdframed}[backgroundcolor=blue!5]
\textbf{Chapter 42} quantizes the photon field in Coulomb (radiation/transverse) gauge.
This gauge makes the physical degrees of freedom manifest --- only two transverse
polarizations propagate --- but breaks manifest Lorentz covariance.  The Coulomb gauge
is the natural precursor to spinor-helicity methods (Ch.~48): helicity polarization
vectors are most cleanly defined in the transverse-gauge framework, and the
Coulomb photon propagator directly exhibits the projector
$\delta^{ij} - \hat p^i\hat p^j$ that underlies the $\langle ij\rangle/[ij]$ brackets.
\end{mdframed}

%% ============================================================
\section{Coulomb Gauge Condition}
%% ============================================================

\textit{Srednicki eq.\ (55.1)--(55.5).}

The Maxwell Lagrangian with an external current $J^\mu$:
\begin{equation}\label{eq:maxwell-lag}
  \mathcal{L} = -\tfrac{1}{4}F_{\mu\nu}F^{\mu\nu} + J^\mu A_\mu,
  \qquad F_{\mu\nu} = \partial_\mu A_\nu - \partial_\nu A_\mu.
\end{equation}

\subsection{The gauge condition}

\begin{mdframed}[backgroundcolor=blue!5]
\textbf{Coulomb (radiation) gauge:}
\begin{equation}\label{eq:coulomb-cond}
  \nabla\cdot\mathbf{A} = \partial_i A^i = 0.
\end{equation}
\end{mdframed}

This gauge condition eliminates the longitudinal photon degree of freedom and
retains only the two physical transverse modes.  The temporal component $A^0$ is
not a dynamical variable --- it is determined by the constraint equation:
\begin{equation}\label{eq:coulomb-A0}
  -\nabla^2 A^0 = J^0 = \rho \qquad \Rightarrow \qquad
  A^0(\mathbf{x},t) = \int d^3y\,\frac{\rho(\mathbf{y},t)}{4\pi|\mathbf{x}-\mathbf{y}|}.
\end{equation}

\subsection{Momentum-space projection}

In momentum space, the Coulomb gauge condition \eqref{eq:coulomb-cond} becomes
$\mathbf{k}\cdot\tilde{\mathbf{A}}(\mathbf{k}) = 0$, implemented by the
\emph{transverse projector}:

\begin{mdframed}[backgroundcolor=blue!5]
\begin{equation}\label{eq:transverse-projector}
  T^{ij}(\mathbf{k}) = \delta^{ij} - \frac{k^i k^j}{|\mathbf{k}|^2}.
\end{equation}
\end{mdframed}

This projector satisfies:
\begin{align}
  k_i T^{ij} &= 0 \quad\text{(transversality)}, \\
  T^{ij}T^{jk} &= T^{ik} \quad\text{(idempotent)}, \\
  T^{ii} &= 2 \quad\text{(two transverse DOF)}.
\end{align}

\begin{mdframed}[backgroundcolor=green!5]
\verified{NumPy:} For $\mathbf{k}=(1,0,0)$: $T^{ij} = \mathrm{diag}(0,1,1)$.
$k_i T^{ij}=0$ \checkmark, $T^{ii}=2$ \checkmark, $T^2=T$ \checkmark.
For general $\mathbf{k}=(k_1,k_2,k_3)$ all properties verified. \checkmark
\end{mdframed}

%% ============================================================
\section{Equal-Time Commutators}
%% ============================================================

\textit{Srednicki eq.\ (55.18)--(55.23).}

The canonical momentum for the vector field is:
\begin{equation}
  \Pi_i = \frac{\partial\mathcal{L}}{\partial\dot A_i} = \dot A_i = E^i \quad
  \text{(transverse part)}.
\end{equation}

\begin{mdframed}[backgroundcolor=blue!5]
\textbf{Equal-time commutators in Coulomb gauge:}
\begin{align}
  [A_i(\mathbf{x},t),\, A_j(\mathbf{y},t)] &= 0, \label{eq:comm-AA}\\
  [\Pi_i(\mathbf{x},t),\, \Pi_j(\mathbf{y},t)] &= 0, \label{eq:comm-PP}\\
  [A_i(\mathbf{x},t),\, \Pi_j(\mathbf{y},t)] &= i\int\!\frac{d^3k}{(2\pi)^3}
  e^{i\mathbf{k}\cdot(\mathbf{x}-\mathbf{y})}
  \left(\delta_{ij} - \frac{k_ik_j}{|\mathbf{k}|^2}\right). \label{eq:comm-AP}
\end{align}
\end{mdframed}

The modified commutator \eqref{eq:comm-AP} (instead of $i\delta_{ij}\delta^3(\mathbf{x}-\mathbf{y})$)
reflects the Coulomb-gauge constraint: only transverse modes are quantized.

\subsection{Mode expansion}

The transverse vector field mode expansion:
\begin{equation}\label{eq:mode-expansion}
  \mathbf{A}(x) = \sum_{\lambda=\pm}\int\!\widetilde{dk}\,
  \bigl[\boldsymbol\varepsilon^*_\lambda(\mathbf{k})\,a_\lambda(\mathbf{k})\,e^{ikx}
        + \boldsymbol\varepsilon_\lambda(\mathbf{k})\,a^\dagger_\lambda(\mathbf{k})\,e^{-ikx}\bigr],
\end{equation}
where $k^0 = |\mathbf{k}|$ (massless), $\widetilde{dk} = \frac{d^3k}{(2\pi)^3 2k^0}$,
and the helicity polarization vectors satisfy:
\begin{equation}\label{eq:helicity-pol}
  \boldsymbol\varepsilon_\pm(\mathbf{k}) = \mp\frac{1}{\sqrt{2}}(\hat e_1 \pm i\hat e_2),
  \qquad
  \mathbf{k}\cdot\boldsymbol\varepsilon_\pm = 0, \qquad
  \boldsymbol\varepsilon_+\cdot\boldsymbol\varepsilon_+^* = 1.
\end{equation}

Polarization completeness (sum over two physical polarizations):
\begin{equation}\label{eq:pol-completeness}
  \sum_{\lambda=\pm}\varepsilon_{i,\lambda}(\mathbf{k})\varepsilon_{j,\lambda}^*(\mathbf{k})
  = \delta_{ij} - \frac{k_ik_j}{|\mathbf{k}|^2} = T^{ij}(\mathbf{k}).
\end{equation}

%% ============================================================
\section{The Photon Propagator in Coulomb Gauge}
%% ============================================================

\textit{Srednicki eq.\ (55.24)--(55.31).}

\subsection{Components}

\begin{mdframed}[backgroundcolor=blue!5]
\textbf{Photon propagator in Coulomb gauge:}
\begin{align}
  D^{00}(k) &= \frac{-1}{|\mathbf{k}|^2}, \label{eq:D00}\\[4pt]
  D^{0i}(k) &= D^{i0}(k) = 0, \label{eq:D0i}\\[4pt]
  D^{ij}(k) &= \frac{1}{k^2+i\varepsilon}\left(\delta^{ij}
              - \frac{k^ik^j}{|\mathbf{k}|^2}\right)
              = \frac{T^{ij}(\mathbf{k})}{k^2+i\varepsilon}. \label{eq:Dij}
\end{align}
\end{mdframed}

\subsection{Physical interpretation}

The propagator has two distinct contributions:

\begin{enumerate}
\item \textbf{Coulomb component} $D^{00}$: This is the \emph{instantaneous} Coulomb
  interaction.  It is \emph{independent of} $k^0$, meaning it acts at equal times ---
  it is not a propagating mode.  This is the familiar $1/|\mathbf{k}|^2$ Coulomb potential
  in momentum space.

\item \textbf{Transverse component} $D^{ij}$: These are the two physical photon
  propagation modes.  The transverse projector $T^{ij}$ ensures only physical
  polarizations propagate.
\end{enumerate}

\subsection{Comparison with Feynman gauge}

In Feynman gauge $(\partial^\mu A_\mu = 0)$, the photon propagator is:
\begin{equation}\label{eq:feynman-gauge-prop}
  D_F^{\mu\nu}(k) = \frac{-\eta^{\mu\nu}}{k^2+i\varepsilon}.
\end{equation}

The Coulomb gauge propagator can be related to the Feynman gauge propagator by:
\begin{equation}
  D_C^{\mu\nu}(k) = D_F^{\mu\nu}(k) + \text{(gauge terms)},
\end{equation}
where the gauge terms vanish for conserved currents $J^\mu$ satisfying $k_\mu J^\mu=0$
(Ward identity).  Physical amplitudes are gauge-invariant and can be computed in
either gauge.

%% ============================================================
\section{The Instantaneous Coulomb Interaction}
%% ============================================================

\textit{Srednicki eq.\ (55.19), (55.24).}

\subsection{Coulomb Hamiltonian}

The Hamiltonian in Coulomb gauge has an explicit instantaneous piece:

\begin{mdframed}[backgroundcolor=blue!5]
\begin{equation}\label{eq:coulomb-hamiltonian}
  H = \sum_{\lambda=\pm}\int\!\widetilde{dk}\,\omega\, a^\dagger_\lambda(\mathbf{k})a_\lambda(\mathbf{k})
    - \int d^3x\,\mathbf{J}(\mathbf{x})\cdot\mathbf{A}(\mathbf{x})
    + H_\text{Coul},
\end{equation}
where the Coulomb energy is:
\begin{equation}\label{eq:coulomb-energy}
  H_\text{Coul} = \frac{1}{2}\int d^3x\,d^3y\;
  \frac{\rho(\mathbf{x},t)\,\rho(\mathbf{y},t)}{4\pi|\mathbf{x}-\mathbf{y}|}.
\end{equation}
\end{mdframed}

The $H_\text{Coul}$ term is the familiar instantaneous Coulomb interaction between
charge densities.  It is \emph{not} Lorentz invariant by itself; the non-invariance
is compensated by the $D^{00}$ piece of the propagator.  The full S-matrix is Lorentz
invariant despite the non-covariant gauge.

\subsection{Effective Coulomb potential}

In position space, $H_\text{Coul}$ gives the potential between point charges:
\begin{equation}
  V(r) = \frac{e_1 e_2}{4\pi r} \quad (r = |\mathbf{x}-\mathbf{y}|).
\end{equation}
The non-locality $(\sim 1/r)$ is the signature of the \emph{instantaneous} gauge:
the Coulomb potential propagates at infinite speed in this gauge, but the
observable effects are still causal because they combine with retarded transverse
photon exchange.

%% ============================================================
\section{Ward Identities in Coulomb Gauge}
%% ============================================================

\textit{Srednicki eq.\ (55.13)-(55.15).}

\subsection{Transversality of physical polarizations}

The physical polarization vectors in Coulomb gauge satisfy:
\begin{equation}\label{eq:transverse-pol}
  \mathbf{k}\cdot\boldsymbol\varepsilon_\lambda(\mathbf{k}) = 0,
\end{equation}
which is the Coulomb-gauge form of the Ward identity $k^\mu\mathcal{M}_\mu = 0$.

\subsection{Ward identity}

For QED amplitudes with an external photon of momentum $k$ and polarization
$\varepsilon^\mu$, current conservation $\partial_\mu J^\mu = 0$ gives:
\begin{equation}\label{eq:ward}
  k_\mu \mathcal{M}^\mu = 0.
\end{equation}
In Coulomb gauge, this means:
\begin{align}
  k_i \mathcal{M}^i &= 0 \quad\text{(transverse)}, \\
  \mathcal{M}^0 &= 0 \quad\text{(temporal component vanishes for on-shell photons)}.
\end{align}

This Ward identity is the foundation for the replacement of $\sum_\lambda \varepsilon_\mu\varepsilon_\nu^*$
with $-\eta_{\mu\nu}$ in Feynman gauge or with $T_{ij}$ in Coulomb gauge.

%% ============================================================
\section{Connection to Light-Front Gauge and Spinor-Helicity}
%% ============================================================

\textit{Bridge from Ch.~42 to Ch.~48.}

\subsection{Light-front gauge}

The Coulomb gauge is related to \emph{light-front} (axial) gauge via the gauge
parameter $n^\mu = (1,0,0,1)$ (the light-front direction):
\begin{equation}
  n^\mu A_\mu = A^+ = 0.
\end{equation}
In light-front gauge, the transverse projector becomes:
\begin{equation}
  T^{\mu\nu}(k;n) = -\eta^{\mu\nu}
  + \frac{k^\mu n^\nu + n^\mu k^\nu}{k\cdot n}
  - \frac{n^2 k^\mu k^\nu}{(k\cdot n)^2}.
\end{equation}
For massless on-shell gluons with $k^2=0$, this simplifies and the helicity
polarization vectors become:

\begin{mdframed}[backgroundcolor=blue!5]
\textbf{Spinor-helicity polarizations} (the Coulomb/light-front gauge origin):
\begin{equation}\label{eq:helicity-pol-spinor}
  \varepsilon^{+\mu}(k;q) = \frac{\langle q|\gamma^\mu|k]}{\sqrt{2}\langle qk\rangle},
  \qquad
  \varepsilon^{-\mu}(k;q) = \frac{[q|\gamma^\mu|k\rangle}{\sqrt{2}[qk]},
\end{equation}
where $q^\mu$ is a reference (gauge) momentum.  These are the spinor-helicity
polarization vectors of Ch.~48, \emph{derived from} the transverse gauge
of Ch.~42.
\end{mdframed}

\subsection{The transverse projector and brackets}

The completeness relation \eqref{eq:pol-completeness}:
\begin{equation}
  \sum_{\lambda=\pm}\varepsilon_{i}^\lambda\varepsilon_j^{*\lambda}
  = \delta_{ij} - \frac{k_ik_j}{|\mathbf{k}|^2} = T_{ij}(\mathbf{k}),
\end{equation}
is the spatial version of the spinor-helicity identity:
\begin{equation}
  \sum_{\lambda=\pm}\varepsilon_\mu^\lambda(k;q)\varepsilon_\nu^{*\lambda}(k;q)
  = -\eta_{\mu\nu} + \frac{k_\mu q_\nu + q_\mu k_\nu}{k\cdot q}.
\end{equation}

In spinor-helicity language, $T_{ij}$ is the projection onto the space spanned by
$(\hat e_+ \equiv \lambda,\, \hat e_- \equiv \tilde\lambda)$ --- exactly the angle
and square spinors of Ch.~48.

\subsection{Summary: Ch.~42 $\to$ Ch.~48}

\begin{center}
\begin{tabular}{ll}
\toprule
Ch.~42 (Coulomb gauge) & Ch.~48 (spinor-helicity) \\
\midrule
Transverse projector $\delta^{ij}-k^ik^j/|\mathbf{k}|^2$ & $\sum_\lambda\varepsilon_\mu\varepsilon_\nu^*$ \\
Coulomb condition $\partial_i A^i=0$ & Helicity gauge $k\cdot\varepsilon=0$ \\
Ward identity $k_i\mathcal{M}^i=0$ & $k\cdot\mathcal{M}=0$ \\
Helicity pol.\ $\varepsilon_\pm = \frac{1}{\sqrt2}(\hat e_1\pm i\hat e_2)$ & $\varepsilon^\pm_\mu=\langle q|\gamma_\mu|k]/(\sqrt2\langle qk\rangle)$ \\
Photon propagator $T^{ij}/k^2$ & Gluon propagator $-\eta^{\mu\nu}/k^2$ (Feynman) \\
\bottomrule
\end{tabular}
\end{center}

\begin{mdframed}[backgroundcolor=blue!5]
\textbf{Key conceptual link:} The spinor-helicity brackets $\langle jk\rangle$ and $[jk]$
of Ch.~48 are the momentum-space representatives of the transverse polarization degrees
of freedom revealed by the Coulomb gauge of Ch.~42.  The rank-1 factorization
$p_{\alpha\dot\alpha}=\lambda_\alpha\tilde\lambda_{\dot\alpha}$ is the covariantization
of the transverse decomposition $\mathbf{A}_\perp = $ (two polarizations).
\end{mdframed}

\vfill
\bibliographystyle{plain}
\end{document}
